% 第1章贯穿图：V={1,2,3,4}，E={{1,2},{1,3},{2,3},{2,4},{3,4}}的三种画法。
\begin{tikzpicture}[>=Stealth,
  v/.style={circle,draw=FigureStroke,fill=FigureBlueFill,line width=.8pt,minimum size=6.5mm,inner sep=0pt},
  e/.style={draw=FigureStroke,line width=.8pt},
  note/.style={font=\FigureNoteFont\fontsize{7.5}{9}\selectfont,text=FigureMuted}]
\begin{scope}
  \node[v] (1) at (-1.3,0) {$1$}; \node[v] (2) at (0,.9) {$2$};
  \node[v] (3) at (0,-.9) {$3$}; \node[v] (4) at (1.3,0) {$4$};
  \draw[e] (1)--(2) (1)--(3) (2)--(3) (2)--(4) (3)--(4);
  \node at (0,-1.75) {(a) 菱形画法};
\end{scope}
\begin{scope}[xshift=3.9cm]
  \node[v] (1) at (-1.35,0) {$1$}; \node[v] (2) at (-.45,0) {$2$};
  \node[v] (3) at (.45,0) {$3$}; \node[v] (4) at (1.35,0) {$4$};
  \draw[e] (1)--(2) (2)--(3) (3)--(4);
  \draw[e] (1) to[bend left=55] (3);
  \draw[e] (2) to[bend right=55] (4);
  \node at (0,-1.75) {(b) 排成一行};
\end{scope}
\begin{scope}[xshift=7.6cm]
  \node[v] (1) at (-.8,.8) {$1$}; \node[v] (4) at (.8,.8) {$4$};
  \node[v] (3) at (-.8,-.8) {$3$}; \node[v] (2) at (.8,-.8) {$2$};
  \draw[e] (1)--(2) (1)--(3) (2)--(3) (2)--(4) (3)--(4);
  \node at (0,-1.75) {(c) 边在画面上交叉};
\end{scope}
\end{tikzpicture}
